%----------------------------------------------------------------------------------------
% HOOFDSTELLING EN PRIMITIEVE FUNCTIES. Alle opgaven zijn eigen werk.
%
% Opbouw: Start, leerstof, en op een nieuwe pagina de Opdrachten, Koppel en Test jezelf.
% Labels van dit hoofdstuk beginnen met h2-.
%----------------------------------------------------------------------------------------

\mychapter{2}{Hoofdstelling en primitieve functies}

\lkr{Leerplan B+S'': LPD 32, de hoofdstelling van de integraalrekening. Het leerplan noemt ze uitdrukkelijk als voorbeeld van een bewijs waarbij beargumenteren belangrijker is dan reproduceren: laat de leerlingen de stappen verantwoorden, of het bewijs herhalen met andere letters. De rekentechnieken beperken zich hier tot de basisprimitieven en de lineariteit; substitutie en partiële integratie volgen in hoofdstuk 3. De oppervlakte tussen krommen staat hier, omdat ze meteen oefent op wat de integraal betekent. Oneigenlijke integralen staan in hoofdstuk 4.}

In hoofdstuk 1 kostte de oppervlakte onder $x^2$ een halve pagina rekenwerk, en voor $\sqrt{x}$ of $\sin x$ lukt die aanpak zelfs niet meer. Dit hoofdstuk lost dat op met een van de mooiste stellingen uit de wiskunde.

Oppervlakte berekenen en raaklijnen zoeken lijken niets met elkaar te maken. Toch ontdekten Newton en Leibniz in de 17de eeuw dat integreren en afleiden elkaars omgekeerde zijn. Daardoor wordt een integraal berekenen plots iets anders: je zoekt een functie waarvan de afgeleide gegeven is. Met die techniek berekenen we op het einde van het hoofdstuk oppervlaktes die met de meetkunde onbereikbaar waren.

\begin{oefeningen}[Start]
\oef{h2-start-1}

\oef{h2-verken}
\end{oefeningen}

\lkr{Oefening \ref{oef:h2-verken} bevat de hele stelling: vraag c is deel 1, vraag d is deel 2. Laat de leerlingen vermoeden of dat toeval is bij deze ene functie, en geef dan pas de stelling.}

%========================================================================================
\section{De hoofdstelling}

\subsection{De integraalfunctie}

Vervang de bovengrens door een veranderlijke $x$: dan hangt de waarde van de integraal af van $x$, en heb je een nieuwe functie. In de integrand staat dan een andere letter, $t$, omdat de $x$ al bezet is.

Hoe snel groeit die functie? Schuif de bovengrens een klein stukje $h$ op. Dan komt er een smalle strook bij, met breedte $h$ en een hoogte van ongeveer $f(x)$, dus de oppervlakte neemt toe met ongeveer $f(x) \cdot h$. De toename per eenheid is dus $f(x)$. Dat is het idee achter de stelling; het bewijs maakt het nauwkeurig.

\needspace{16\baselineskip}
\begin{theorie}
Is $f$ continu op $[a, b]$, dan is de \textbf{\textcolor{accent}{integraalfunctie}} van $f$: \quad $\formule{A(x) = \int_a^x f(t)\dt}$

\begin{stelling}{Hoofdstelling, deel 1}
\begin{minipage}[t]{.62\linewidth}
\vspace{0pt}
De afgeleide van de integraalfunctie is de functie zelf:

\vspace{1mm}
\qquad $\formule{A'(x) = f(x)}$ \quad voor elke $x \in [a, b]$

\vspace{1.5mm}
of, in de notatie van Leibniz:

\vspace{1mm}
\qquad $\formule{\dfrac{\mathrm{d}}{\mathrm{d}x} \int_a^x f(t)\dt = f(x)}$
\end{minipage}\hfill
\begin{minipage}[t]{.34\linewidth}
\vspace{0pt}
\begin{tikzpicture}[x=1cm,y=.62cm,font=\scriptsize]
  \gebied[accent!18]{.5}{2}{1+.5*\x+.15*(\x)^2}
  \gebied[accent!45]{2}{2.5}{1+.5*\x+.15*(\x)^2}
  \assen{-.1}{3.8}{-.1}{4.2}
  \kromme{.2}{3.4}{1+.5*\x+.15*(\x)^2}
  \node[below] at (.5,0) {$a$};
  \node[below] at (2,0) {$x$};
  \node[below,inner xsep=0pt] at (2.6,0) {$x + h$};
  \node at (1.3,.8) {$A(x)$};
\end{tikzpicture}
\end{minipage}
\end{stelling}

\begin{bewijs}
Neem $x,\ x + h \in [a, b]$ met $h \neq 0$. Wegens de optelbaarheid is de toename van $A$ de donkere strook in de figuur:
\bewijsregel{A(x + h) - A(x) = \ig_a^{x+h} f(t)\dt - \ig_a^{x} f(t)\dt = \ig_x^{x+h} f(t)\dt .}
Is $h > 0$, dan geeft de middelwaardestelling op $[x, x + h]$ een $c$ met
\bewijsregel{\ig_x^{x+h} f(t)\dt = f(c) \cdot h .}
Is $h < 0$, dan ligt $x + h$ links van $x$. De middelwaardestelling op $[x + h, x]$, met lengte $-h$, geeft dan een $c$ met
\bewijsregel{\ig_x^{x+h} f(t)\dt = -\ig_{x+h}^{x} f(t)\dt = -f(c) \cdot (-h) = f(c) \cdot h .}
In beide gevallen ligt $c$ tussen $x$ en $x + h$, en is
\bewijsregel{\dfrac{A(x + h) - A(x)}{h} = f(c) .}
Laat nu $h \to 0$. Dan $c \to x$, en omdat $f$ continu is, $f(c) \to f(x)$. Bijgevolg is
\bewijsregel{A'(x) = \lim_{h \to 0} \dfrac{A(x + h) - A(x)}{h} = \lim_{h \to 0} f(c) = f(x) .}
\end{bewijs}
\end{theorie}

Wat zegt deze stelling nu eigenlijk? De verandering van de oppervlakte is precies gelijk aan de functiewaarde. Wat er bijkomt als $x$ een klein beetje opschuift, kun je zien als een infinitesimale strook met hoogte $f(x)$. Hoe hoger de grafiek daar ligt, hoe sneller de oppervlakte groeit. Informeel gezegd: integreren is het omgekeerde van afleiden. Integreer je $f$ en leid je het resultaat af, dan krijg je $f$ terug.

De stelling zegt ook hoe de grafieken van $f$ en $A$ samenhangen. In de figuur is $f(t) = 2 - t$ en $A(x) = \int_0^x f(t)\dt$. Zolang $f$ positief is, komt er oppervlakte bij en stijgt $A$. Waar $f$ nul wordt, bereikt $A$ zijn top. Daarna telt er negatieve oppervlakte mee en daalt $A$. De functie $f$ is overal de helling van $A$.

\figuren{\begin{tikzpicture}[x=.75cm,y=.55cm,font=\scriptsize]
  \gebied{0}{2}{2-\x}
  \gebied[accent!15]{2}{4}{2-\x}
  \rooster{0}{1}{4}{-2}{1}{2}
  \kromme{0}{4}{2-\x}
  \node[above,text=accent] at (2.8,.9) {$y = f(t)$};
\end{tikzpicture}\qquad\qquad
\begin{tikzpicture}[x=.75cm,y=.55cm,font=\scriptsize]
  \rooster{0}{1}{4}{-2}{1}{2}
  \kromme{0}{4}{2*\x-.5*(\x)^2}
  \node[above,text=accent] at (2,2.1) {$y = A(x)$};
  \draw[lijngrijs,dashed] (2,-2) -- (2,2);
\end{tikzpicture}}

Het mooie is dat je zo kunt afleiden zonder de integraal zelf te kennen.

\voorbeeld{}{Bereken de afgeleide van $A(x) = \ig_1^x \sqrt{1 + t^3}\dt$ en van $B(x) = \ig_1^{x^2} \sqrt{1 + t^3}\dt$.\label{vb:h2-afgeleide}}{%
\uitleg{De integraal zelf kunnen we niet uitrekenen, en dat hoeft ook niet: de hoofdstelling geeft de afgeleide rechtstreeks. Let wel op of de bovengrens gewoon $x$ is.}
\begin{stappen}
A'(x) = \sqrt{1 + x^3} & De bovengrens is $x$, dus de hoofdstelling geldt meteen: vervang $t$ in de integrand door $x$. \\[1ex]
B(x) = A(x^2) & Hier is de bovengrens $x^2$. $B$ is dus de functie $A$ van daarnet, ingevuld in $x^2$: een samengestelde functie. \\[1ex]
B'(x) = A'(x^2) \cdot 2x = \sqrt{1 + x^6} \cdot 2x & Daarom de kettingregel: $A'$ ingevuld in $x^2$, maal de afgeleide van $x^2$.
\end{stappen}}

In zo'n integraal staan twee soorten letters, en die moet je goed uit elkaar houden. De grensvariabele $x$ staat in de grens: daarvan hangt de uitkomst af, en daarnaar leid je af. De dummyvariabele $t$ bestaat alleen binnen de integraal: ze loopt van de ondergrens naar de bovengrens, en is na het integreren weg. Daarom mag je haar door elke andere letter vervangen: $\int_1^x \sin(t^2)\dt = \int_1^x \sin(u^2)\du$.

\voorbeeld{}{Bereken de afgeleide van $F(x) = \ig_1^x x \sin(t^2)\dt$.\label{vb:h2-letters}}{%
\uitleg{Hier staat de grensvariabele $x$ ook in de integrand, en dan mag je niet zomaar $t$ door $x$ vervangen. Voor de integraal is $t$ de veranderlijke, en is $x$ een constante.}
\begin{stappen}
F(x) = x \ig_1^x \sin(t^2)\dt & Tijdens het integreren verandert alleen $t$, dus $x$ is een constante factor, en die komt voor de integraal. \\[2.5ex]
F'(x) = 1 \cdot \ig_1^x \sin(t^2)\dt + x \cdot \sin(x^2) & Nu is $F$ een product van twee functies van $x$, dus we leiden af met de productregel. De hoofdstelling geeft de afgeleide van de integraal.
\end{stappen}
Het blinde antwoord $F'(x) = x\sin(x^2)$ mist dus een term. En let op: $\int_1^3 \sin(t^2)\dt$ bevat geen $x$ meer. Dat is een getal, met afgeleide 0.}

\nujij{Bereken de afgeleide van $G(x) = \int_x^{5} \sin(t^2)\dt$.}{Grenzen wisselen: $G(x) = -\int_5^x \sin(t^2)\dt$, dus $G'(x) = -\sin(x^2)$.}

\begin{oefeningen}
\oef{h2-afgeleide}
\end{oefeningen}

\subsection{Primitieve functies}

Deel 1 van de hoofdstelling levert een functie op waarvan de afgeleide $f$ is. Hoeveel zulke functies zijn er eigenlijk? Onderzoek dat eerst zelf.

\begin{oefeningen}
\oef{h2-prim-verken}
\end{oefeningen}

\naast[4.2cm]{Zo'n functie krijgt een eigen naam: een primitieve. Er is er dus nooit maar één. Zowel $x^2$ als $x^2 + 1$ en $x^2 - 2$ hebben als afgeleide $2x$. Hun grafieken zijn dezelfde parabool, verticaal verschoven: boven eenzelfde $x$ hebben ze overal dezelfde helling. En twee zulke functies verschillen altijd hoogstens een constante: meer primitieven zijn er niet.

Hoe helpt dat om een integraal te berekenen? De integraalfunctie is een van die primitieven. We hoeven dus maar één primitieve te vinden, en de constante te bepalen.}{\begin{tikzpicture}[x=.75cm,y=.42cm,font=\scriptsize]
  \rooster{-2}{1}{2}{-2}{1}{5}
  \foreach \c in {-2,-1,0,1}{\kromme{-1.95}{1.95}{(\x)^2+\c}}
  \kromme[very thick]{-1.95}{1.95}{(\x)^2}
  \foreach \c in {-2,-1,0,1}{\draw[grijs] (.6,{1+\c-.8}) -- (1.4,{1+\c+.8});}
\end{tikzpicture}}


\begin{theorie}
$F$ is een \textbf{\textcolor{accent}{primitieve functie}} van $f$ op een interval als $\formule{F'(x) = f(x)}$ op dat interval.

\begin{stelling}{Verschil van twee primitieven}
Twee primitieven van dezelfde functie op een interval verschillen hoogstens een constante: \quad $\formule{G(x) = F(x) + C}$
\end{stelling}

\begin{bewijs}
Stel $F' = G' = f$ en $H = G - F$. Dan is $H' = f - f = 0$ op het interval. Een functie met afgeleide 0 op een interval is constant, zeg $H(x) = C$, en daaruit volgt $G(x) = F(x) + C$.
\end{bewijs}

\begin{stelling}{Hoofdstelling, deel 2}
Is $f$ continu op $[a, b]$ en $F$ een primitieve van $f$, dan

\vspace{1mm}
\qquad $\formule{\int_a^b f(x)\dx = \Bigl[F(x)\Bigr]_a^b = F(b) - F(a)}$
\end{stelling}

\begin{bewijs}
Volgens deel 1 is $A$ een primitieve van $f$, dus $A(x) = F(x) + C$. Voor $x = a$ geeft dat
\bewijsregel{0 = \ig_a^a f(t)\dt = A(a) = F(a) + C \quad\Rightarrow\quad C = -F(a) ,}
en voor $x = b$ krijgen we dan
\bewijsregel{\ig_a^b f(t)\dt = A(b) = F(b) - F(a) .}
\end{bewijs}
\end{theorie}

\voorbeeld{}{Welke van deze functies zijn een primitieve van $f(x) = 3x^2 + 1$: \ $x^3 + x$, \ $x^3 + x - 5$, \ $6x$?\label{vb:h2-prim}}{%
\uitleg{Om te controleren of iets een primitieve is, leid je het af. Krijg je $f$ terug, dan is het een primitieve.}
\begin{stappen}
(x^3 + x)' = 3x^2 + 1 & Dat is $f(x)$, dus dit is een primitieve. \\[1ex]
(x^3 + x - 5)' = 3x^2 + 1 & Ook een primitieve. Ze verschilt van de vorige de constante $-5$, en die valt weg bij het afleiden. \\[1ex]
(6x)' = 6 & Dat is niet $f(x)$, dus dit is geen primitieve. Omgekeerd is $6x$ de afgeleide van $f$: je bent de verkeerde kant op gegaan.
\end{stappen}}

Die stelling is bijzonder: de limiet van een som, waar je in hoofdstuk 1 een hele pagina voor nodig had, is nu een verschil van twee functiewaarden. De constante $C$ doet er niet toe, want ze valt weg bij het aftrekken. Daarom kies je $C = 0$.

\voorbeeld{}{Bereken $\ig_0^3 x^2\dx$. Vergelijk met voorbeeld 2 uit hoofdstuk 1.\label{vb:h2-hoofd}}{%
\uitleg{We zoeken een functie met afgeleide $x^2$. Bij een macht stijgt de exponent met 1, en delen we door de nieuwe exponent.}
\naast[3.6cm]{\begin{stappenvrij}{6.8cm}{4cm}
F(x) = \dfrac{x^3}{3} & Controleer altijd door af te leiden: $F'(x) = x^2$. \\[2ex]
\ig_0^3 x^2\dx = \grens{\dfrac{x^3}{3}}{0}{3} = \dfrac{27}{3} - 0 = 9 & Bovengrens min ondergrens. Dezelfde uitkomst als in hoofdstuk 1, maar in één regel.
\end{stappenvrij}}{\begin{tikzpicture}[x=.75cm,y=.3cm,font=\scriptsize]
  \gebied{0}{3}{(\x)^2}
  \rooster{0}{1}{3}{0}{3}{9}
  \kromme{0}{3}{(\x)^2}
\end{tikzpicture}}}

Tel de hokjes in de figuur na: 9 is een geloofwaardige waarde. Die controle kost weinig tijd, en ze vangt veel rekenfouten op.

\begin{werkwijze}
Zoek een primitieve $F$. Schrijf ze tussen rechte haken met de grenzen erbij. Vul eerst de bovengrens en dan de ondergrens in, en trek af.
\end{werkwijze}

\nujij{Bereken $\int_1^2 x^3\dx$.}{$\bigl[\tfrac{x^4}{4}\bigr]_1^2 = 4 - \tfrac14 = \tfrac{15}{4}$.}

\begin{oefeningen}
\oef{h2-hoofd}
\end{oefeningen}

%========================================================================================
\section{Integralen berekenen}

\subsection{Basisprimitieven}

Om de hoofdstelling te gebruiken, heb je primitieven nodig. Die zoek je niet met een nieuwe methode: elke afgeleide die je kent, geeft een primitieve als je ze omgekeerd leest. Omdat $(x^3)' = 3x^2$, is $\tfrac{x^3}{3}$ een primitieve van $x^2$.

\begin{theorie}
\begin{stelling}{Basisprimitieven}
{\centering\arrayrulecolor{accent}\renewcommand{\arraystretch}{2.7}%
\begin{tabular}{|>{\centering\arraybackslash$}m{3.3cm}<{$}|>{\centering\arraybackslash$}m{2.3cm}<{$}|>{\centering\arraybackslash$}m{2.8cm}<{$}|>{\centering\arraybackslash$}m{2.3cm}<{$}|}
\hline
\mathbf{f(x)} & \mathbf{F(x)} & \mathbf{f(x)} & \mathbf{F(x)} \\ \hline
\formule{x^r} \ (r \neq -1) & \formule{\dfrac{x^{r+1}}{r+1}} & \formule{\sin x} & \formule{-\cos x} \\ \hline
\formule{\dfrac{1}{x}} & \formule{\ln |x|} & \formule{\cos x} & \formule{\sin x} \\ \hline
\formule{\e^x} & \formule{\e^x} & \formule{\dfrac{1}{\cos^2 x}} & \formule{\tan x} \\ \hline
\formule{a^x} \ (a > 0,\ a \neq 1) & \formule{\dfrac{a^x}{\ln a}} & \formule{\dfrac{1}{\sin^2 x}} & \formule{-\cot x} \\ \hline
\formule{\dfrac{1}{1 + x^2}} & \formule{\bgtan x} & \formule{\dfrac{1}{\sqrt{1 - x^2}}} & \formule{\bgsin x} \\ \hline
\end{tabular}\par}
\end{stelling}
\end{theorie}

\naast[4.6cm]{Twee dingen vragen wat uitleg. Waarom staat er $\ln |x|$? De functie $\tfrac1x$ bestaat ook voor $x < 0$, en voor die $x$ is $\bigl(\ln(-x)\bigr)' = \tfrac{1}{-x} \cdot (-1) = \tfrac1x$. De figuur toont het: links stijgt $\ln|x|$ niet, maar daalt ze, en daar is $\tfrac1x$ dan ook negatief.

En waar is de macht $r = -1$ gebleven? De eerste regel van de tabel zou delen door 0. Precies die ene macht heeft een logaritme als primitieve.

Een wortel of een breuk schrijf je eerst als een macht, want dan past ze in de eerste regel.

Waarom staat $\bgcos x$ niet in de tabel? Omdat $\bigl(\bgcos x\bigr)' = -\tfrac{1}{\sqrt{1 - x^2}}$, is ze een primitieve van die functie. Maar toevoegen heeft geen zin: ook $-\bgsin x$ is een primitieve ervan, en die heb je al. De twee verschillen een constante, want $\bgcos x = \tfrac{\pi}{2} - \bgsin x$.}{\begin{tikzpicture}[x=.6cm,y=.5cm,font=\scriptsize]
  \rooster{-3}{1}{3}{-3}{1}{3}
  \kromme{.34}{3}{1/\x}
  \kromme{-3}{-.34}{1/\x}
  \kromme[dashed]{.06}{3}{ln(\x)}
  \kromme[dashed]{-3}{-.06}{ln(-\x)}
  \node[text=accent] at (1.1,2.3) {$\frac1x$};
  \node[text=accent] at (2.2,1.45) {$\ln|x|$};
\end{tikzpicture}}

\voorbeeld{}{Bereken $\ig_1^4 \sqrt{x}\dx$, \ $\ig_1^{\e} \dfrac{1}{x}\dx$ en $\ig_0^{1} \dfrac{1}{1 + x^2}\dx$.\label{vb:h2-basis}}{%
\uitleg{Drie integralen die rechtstreeks uit de tabel komen, soms na het herschrijven als macht. Bij een exacte uitkomst laat je $\pi$ of $\e$ staan.}
\begin{stappen}
\ig_1^4 x^{1/2}\dx = \grens{\dfrac{x^{3/2}}{3/2}}{1}{4} = \dfrac23 (8 - 1) = \dfrac{14}{3} & De wortel schrijven we eerst als macht, dan past ze in de eerste regel van de tabel. $4^{3/2} = (\sqrt{4})^3 = 8$. \\[3ex]
\ig_1^{\e} \dfrac{1}{x}\dx = \Bigl[\ln |x|\Bigr]_1^{\e} = 1 - 0 = 1 & Rechtstreeks uit de tabel, met $\ln \e = 1$ en $\ln 1 = 0$. \\[3ex]
\ig_0^{1} \dfrac{1}{1 + x^2}\dx = \Bigl[\bgtan x\Bigr]_0^1 = \dfrac{\pi}{4} & $\bgtan 1$ is de hoek met tangens 1, in radialen: $\tfrac{\pi}{4}$. En $\bgtan 0 = 0$.
\end{stappen}

\vspace{1mm}
\figuren{\begin{tikzpicture}[x=.55cm,y=.55cm,font=\scriptsize]
  \gebied{1}{4}{sqrt(\x)}
  \rooster{0}{1}{4}{0}{1}{2}
  \kromme{0.02}{4}{sqrt(\x)}
\end{tikzpicture}\quad
\begin{tikzpicture}[x=.8cm,y=.8cm,font=\scriptsize]
  \gebied{1}{2.718}{1/\x}
  \rooster{0}{1}{3}{0}{1}{2}
  \kromme{.5}{3}{1/\x}
  \node[below] at (2.718,0) {$\e$};
\end{tikzpicture}\quad
\begin{tikzpicture}[x=1.3cm,y=1.3cm,font=\scriptsize]
  \gebied{0}{1}{1/(1+(\x)^2)}
  \rooster{0}{1}{2}{0}{1}{1}
  \kromme{0}{2}{1/(1+(\x)^2)}
\end{tikzpicture}}
De figuren maken de uitkomsten aannemelijk. Het eerste gebied is een bijna-trapezium met oppervlakte iets meer dan $\tfrac{1 + 2}{2} \cdot 3 = 4{,}5$. Het laatste is iets kleiner dan 1: $\tfrac{\pi}{4} \approx 0{,}79$.}

\nujij{Bereken $\int_0^{\pi/2} \cos x\dx$ en $\int_0^2 3^x\dx$.}{$\sin\tfrac{\pi}{2} - \sin 0 = 1$ en $\bigl[\tfrac{3^x}{\ln 3}\bigr]_0^2 = \tfrac{8}{\ln 3}$.}

\voorbeeld{}{Voor elke regel van de tabel een integraal. (Niet uitgewerkt, reken ze zelf na.)\label{vb:h2-tabel}}{%
\uitleg{Zoek per integraal de juiste regel in de tabel. Let bij de derde op de negatieve grenzen, en bij de vijfde op het grondtal.}
\begin{tabular}{@{}*{2}{>{$\displaystyle}l<{$}@{\hspace{10mm}}}@{}}
\ig_0^1 x^4\dx = \dfrac15 & \ig_0^{\pi/2} \sin x\dx = 1 \\[2.5ex]
\ig_1^4 x^{-1/2}\dx = 2 & \ig_0^{\pi/2} \cos x\dx = 1 \\[2.5ex]
\ig_{-\e}^{-1} \dfrac1x\dx = -1 & \ig_0^{\pi/4} \dfrac{1}{\cos^2 x}\dx = 1 \\[2.5ex]
\ig_0^1 \e^x\dx = \e - 1 & \ig_{\pi/4}^{\pi/2} \dfrac{1}{\sin^2 x}\dx = 1 \\[2.5ex]
\ig_0^2 2^x\dx = \dfrac{3}{\ln 2} & \ig_0^{1/2} \dfrac{1}{\sqrt{1 - x^2}}\dx = \dfrac{\pi}{6}
\end{tabular}}

\begin{oefeningen}
\oef{h2-basis}
\end{oefeningen}

\subsection{Meerdere termen}

Meestal staat er meer dan één term. Dankzij de lineariteit uit hoofdstuk 1 mag je dan term per term een primitieve zoeken. Soms moet je de integrand eerst herschrijven, want er bestaat geen eenvoudige regel voor een product of een quotiënt.

\voorbeeld{}{Bereken $\ig_1^2 \Bigl(3x^2 - \dfrac{2}{x} + \e^x\Bigr)\dx$.\label{vb:h2-lin}}{%
\uitleg{Elke term staat in de tabel, dus de lineariteit doet het werk. De valkuil zit bij het invullen van de ondergrens.}
\naast[3.5cm]{\begin{stappenvrij}{6.8cm}{4.1cm}
\grens{x^3 - 2\ln |x| + \e^x}{1}{2} & Lineariteit: een primitieve per term. \\[2.5ex]
= (8 - 2\ln 2 + \e^2) - (1 - 0 + \e) & Zonder haakjes zou het minteken alleen voor de eerste term staan. \\[1ex]
= 7 - 2\ln 2 + \e^2 - \e & $\ln 1 = 0$. Exact laten staan.
\end{stappenvrij}}{\begin{tikzpicture}[x=1.1cm,y=.13cm,font=\scriptsize]
  \gebied{1}{2}{3*(\x)^2-2/\x+exp(\x)}
  \rooster{0}{1}{2}{0}{5}{20}
  \kromme{.7}{2}{3*(\x)^2-2/\x+exp(\x)}
\end{tikzpicture}}}

Numeriek is die uitkomst ongeveer 10,3: dat past bij de figuur, met een breedte 1 en hoogtes tussen 3,7 en 18,4.

Ook het volgende voorbeeld lijkt moeilijk door de breuk, maar na het uitwerken blijven er enkel machten over.

\voorbeeld{}{Bereken $\ig_1^4 \dfrac{(x + 1)^2}{\sqrt{x}}\dx$.\label{vb:h2-splits}}{%
\uitleg{Een quotiënt heeft geen eigen regel, maar deze noemer bestaat uit één term. Dan kunnen we term per term delen, en blijven er machten over.}
\begin{stappen}
\dfrac{x^2 + 2x + 1}{x^{1/2}} = x^{3/2} + 2x^{1/2} + x^{-1/2} & Werk de teller uit en deel elke term door de noemer. Bij machten trek je de exponenten van elkaar af. \\[2ex]
\grens{\tfrac25 x^{5/2} + \tfrac43 x^{3/2} + 2x^{1/2}}{1}{4} & Per term een primitieve: exponent plus 1, en delen door die nieuwe exponent. \\[2ex]
= \bigl(\tfrac{64}{5} + \tfrac{32}{3} + 4\bigr) - \bigl(\tfrac25 + \tfrac43 + 2\bigr) = \dfrac{356}{15} & Vul in: $4^{1/2} = 2$, $4^{3/2} = 8$ en $4^{5/2} = 32$. Controle met het rekentoestel: ongeveer $23{,}73$.
\end{stappen}}

\begin{werkwijze}
Werk uit tot een som van termen die in de tabel staan: haakjes uitwerken, een breuk splitsen. Zoek per term een primitieve. Zet de ingevulde ondergrens tussen haakjes, want het minteken staat voor alle termen.
\end{werkwijze}

\nujij{Bereken $\int_1^2 \dfrac{x^2 + 1}{x}\dx$.}{Splitsen: $\int_1^2 (x + \tfrac1x)\dx = \bigl[\tfrac{x^2}{2} + \ln |x|\bigr]_1^2 = \tfrac32 + \ln 2$.}

\begin{oefeningen}
\oef{h2-lin}
\end{oefeningen}

\subsection{Wanneer mag het niet?}

\naast[3.6cm]{De hoofdstelling vraagt een functie die continu is op het hele interval. Kijk dus altijd eerst na of de grenzen binnen het domein liggen, en er geen gat tussen zit.

Zo geeft blind invullen $\int_{-1}^{1} \tfrac{1}{x^2}\dx = \bigl[-\tfrac{1}{x}\bigr]_{-1}^{1} = -2$, en dat kan niet. De figuur toont waarom: de functie is overal positief, dus een negatieve integraal is onmogelijk. In $x = 0$ bestaat de functie niet, en dus geldt de stelling niet. Hoofdstuk 4 toont wat je dan wel kunt doen.}{\begin{tikzpicture}[x=.7cm,y=.4cm,font=\scriptsize]
  \gebied{-1}{-.42}{1/(\x)^2}
  \gebied{.42}{1}{1/(\x)^2}
  \rooster{-2}{1}{2}{0}{1}{5}
  \kromme{-2}{-.45}{1/(\x)^2}
  \kromme{.45}{2}{1/(\x)^2}
\end{tikzpicture}}

%========================================================================================
\section{Oppervlakte berekenen}

\subsection{Tussen een grafiek en de horizontale as}

De integraal telt een stuk onder de as negatief. Wil je de echte oppervlakte, dan moet je die stukken omhoog klappen: dat doet de absolute waarde. Van de grafiek links maak je zo de grafiek rechts, en de oppervlakte verandert daarbij niet. Maar $|f|$ heeft geen handige primitieve, dus in de praktijk splits je in de nulpunten.

\figuren{\begin{tikzpicture}[x=.68cm,y=.42cm,font=\scriptsize]
  \gebied{0}{1}{(\x)^2-4*\x+3}
  \gebied[accent!15]{1}{3}{(\x)^2-4*\x+3}
  \gebied{3}{4}{(\x)^2-4*\x+3}
  \rooster{0}{1}{4}{-1}{1}{3}
  \kromme{0}{4}{(\x)^2-4*\x+3}
  \node[text=accent] at (2.4,2.3) {$y = f(x)$};
\end{tikzpicture}\qquad\qquad
\begin{tikzpicture}[x=.68cm,y=.42cm,font=\scriptsize]
  \gebied{0}{1}{(\x)^2-4*\x+3}
  \gebied{1}{3}{-((\x)^2-4*\x+3)}
  \gebied{3}{4}{(\x)^2-4*\x+3}
  \rooster{0}{1}{4}{-1}{1}{3}
  \kromme{0}{1}{(\x)^2-4*\x+3}
  \kromme{1}{3}{-((\x)^2-4*\x+3)}
  \kromme{3}{4}{(\x)^2-4*\x+3}
  \node[text=accent] at (2,2.3) {$y = |f(x)|$};
\end{tikzpicture}}

\begin{theorie}
\begin{stelling}{Oppervlakte tussen grafiek en as}
De oppervlakte tussen de grafiek van $f$ en de horizontale as op $[a, b]$ is \qquad $\formule{A = \int_a^b |f(x)|\dx}$
\end{stelling}
\end{theorie}

\voorbeeld{}{Bereken de oppervlakte tussen de grafiek van $f(x) = x^2 - 4x + 3$ en de horizontale as op $[0, 4]$. Bereken ook de integraal.\label{vb:h2-asopp}}{%
\uitleg{Bij een oppervlakte moet je eerst weten waar de grafiek de as snijdt. Daarna rekenen we elk stuk apart uit, en tellen we de absolute waarden op.}
\begin{minipage}[t]{.70\linewidth}
\begin{tabular}{@{}>{$\displaystyle}l<{$}@{\hspace{3mm}}>{\RaggedRight\small\itshape\arraybackslash}p{4.4cm}@{}}
x^2 - 4x + 3 = 0 \ \Rightarrow \ x = 1,\ x = 3 & De nulpunten splitsen het interval. Positief op $[0, 1]$ en $[3, 4]$, negatief ertussen. \\[.5ex]
F(x) = \dfrac{x^3}{3} - 2x^2 + 3x & Eén primitieve volstaat voor alle stukken. \\[2ex]
\ig_0^1 f = \tfrac43, \quad \ig_1^3 f = -\tfrac43, \quad \ig_3^4 f = \tfrac43 & $F(0) = 0$, $F(1) = \tfrac43$, $F(3) = 0$, $F(4) = \tfrac43$. \\[2.5ex]
A = \tfrac43 + \tfrac43 + \tfrac43 = 4 & Elk stuk telt positief, ook het middelste. \\[1ex]
\ig_0^4 f(x)\dx = \tfrac43 & De integraal houdt het minteken: $\tfrac43 - \tfrac43 + \tfrac43$. Iets heel anders.
\end{tabular}
\end{minipage}\hfill
\begin{minipage}[t]{.26\linewidth}
\vspace{-2mm}
\begin{tikzpicture}[x=.75cm,y=.6cm,font=\scriptsize]
  \gebied{0}{1}{(\x)^2-4*\x+3}
  \gebied[accent!15]{1}{3}{(\x)^2-4*\x+3}
  \gebied{3}{4}{(\x)^2-4*\x+3}
  \assen{-.2}{4.5}{-1.4}{3.6}
  \kromme{0}{4.05}{(\x)^2-4*\x+3}
  \foreach \t in {1,3,4} \node[above] at (\t,0) {\t};
\end{tikzpicture}
\end{minipage}}

\begin{werkwijze}
Zoek de nulpunten van $f$ op $[a, b]$ en maak een tekenschema. Bereken de integraal per stuk tussen twee nulpunten. Maak de stukken positief en tel ze op.
\end{werkwijze}

\nujij{Bereken de oppervlakte tussen de grafiek van $f(x) = x^2 - 1$ en de horizontale as op $[0, 2]$.}{$\int_0^1 f = -\tfrac23$ en $\int_1^2 f = \tfrac43$, dus $A = 2$.}

\begin{oefeningen}
\oef{h2-asopp}
\end{oefeningen}

\subsection{Tussen twee grafieken}

Wat als het gebied niet tot aan de as loopt, maar tussen twee grafieken ligt? Ligt alles boven de as, dan is het de oppervlakte onder de bovenste grafiek min die onder de onderste. En als een deel onder de as ligt? Schuif de twee grafieken samen omhoog tot ze boven de as liggen. De oppervlakte ertussen verandert daarbij niet, en in $f - g$ valt de verschuiving weg. De plaats van de as speelt dus geen rol, en je moet niet splitsen in de nulpunten van $f$ of $g$.

\figuren{\begin{tikzpicture}[x=.7cm,y=.42cm,font=\scriptsize]
  \gebiedtussen{.5}{3.5}{1.2-.3*(\x-2)^2}{.3*(\x-2)^2-1.8}
  \assen{-.1}{4.2}{-2.5}{3.3}
  \kromme{.2}{3.8}{1.2-.3*(\x-2)^2}
  \kromme{.2}{3.8}{.3*(\x-2)^2-1.8}
  \node[above] at (2,1.2) {$f$};
  \node[below] at (2,-1.8) {$g$};
\end{tikzpicture}\qquad\qquad
\begin{tikzpicture}[x=.7cm,y=.42cm,font=\scriptsize]
  \gebiedtussen{.5}{3.5}{3.4-.3*(\x-2)^2}{.3*(\x-2)^2+.4}
  \assen{-.1}{4.2}{-.3}{4.8}
  \kromme{.2}{3.8}{3.4-.3*(\x-2)^2}
  \kromme{.2}{3.8}{.3*(\x-2)^2+.4}
  \node[above] at (2,3.4) {$f + c$};
  \node[below] at (2,.4) {$g + c$};
\end{tikzpicture}}

\begin{theorie}
\begin{stelling}{Oppervlakte tussen twee grafieken}
\begin{minipage}[t]{.62\linewidth}
\vspace{0pt}
De oppervlakte tussen de grafieken van $f$ en $g$ op $[a, b]$ is

\vspace{1mm}
\qquad $\formule{A = \int_a^b |f(x) - g(x)|\dx}$

\vspace{1.5mm}
Ligt $f$ overal boven $g$ op $[a, b]$, dan is dat

\vspace{1mm}
\qquad $\formule{A = \int_a^b \bigl(f(x) - g(x)\bigr)\dx}$
\end{minipage}\hfill
\begin{minipage}[t]{.34\linewidth}
\vspace{0pt}
\begin{tikzpicture}[x=1cm,y=.55cm,font=\scriptsize]
  \gebiedtussen{.6}{3.2}{3.1+.5*sin(\x*90)}{.2*(\x-1.5)^2-.6}
  \assen{-.1}{3.8}{-1}{4.2}
  \kromme{.3}{3.5}{3.1+.5*sin(\x*90)}
  \kromme{.3}{3.5}{.2*(\x-1.5)^2-.6}
  \draw[lijngrijs,dashed] (.6,0) -- (.6,3.5) (3.2,0) -- (3.2,3);
  \node[above left] at (.6,0) {$a$};
  \node[above right] at (3.2,0) {$b$};
  \node[above] at (2.1,3.6) {$f$};
  \node[below] at (2.1,-.5) {$g$};
\end{tikzpicture}
\end{minipage}
\end{stelling}
\end{theorie}

\voorbeeld{}{Bereken de oppervlakte van het gebied ingesloten door de grafieken van $f(x) = x^2$ en $g(x) = x + 2$.\label{vb:h2-tussen}}{%
\uitleg{Er zijn geen grenzen gegeven: de snijpunten bepalen het gebied. Daarna kijken we welke grafiek bovenaan ligt.}
\begin{minipage}[t]{.70\linewidth}
\begin{tabular}{@{}>{$\displaystyle}l<{$}@{\hspace{3mm}}>{\RaggedRight\small\itshape\arraybackslash}p{4.4cm}@{}}
x^2 = x + 2 \ \Rightarrow \ x = -1,\ x = 2 & De snijpunten zijn de grenzen. \\[1ex]
g(0) = 2 > f(0) = 0 & Test een punt tussen de snijpunten, hier $x = 0$. De rechte ligt bovenaan. \\[1.5ex]
A = \ig_{-1}^{2} (x + 2 - x^2)\dx & Bovenste min onderste. \\[2.5ex]
\phantom{A} = \grens{\dfrac{x^2}{2} + 2x - \dfrac{x^3}{3}}{-1}{2} & Per term een primitieve. \\[2.5ex]
\phantom{A} = \dfrac{10}{3} - \Bigl(-\dfrac{7}{6}\Bigr) = \dfrac{9}{2} & Min maal min geeft plus.
\end{tabular}
\end{minipage}\hfill
\begin{minipage}[t]{.26\linewidth}
\vspace{-2mm}
\begin{tikzpicture}[x=.8cm,y=.5cm,font=\scriptsize]
  \gebiedtussen{-1}{2}{\x+2}{(\x)^2}
  \assen{-1.6}{2.6}{-.3}{5}
  \kromme{-1.5}{2.2}{(\x)^2}
  \kromme{-1.5}{2.4}{\x+2}
  \node[below] at (-1,0) {$-1$};
  \node[below] at (2,0) {2};
\end{tikzpicture}
\end{minipage}}

\begin{werkwijze}
Zoek de snijpunten met $f(x) = g(x)$. Bepaal op elk stuk welke grafiek bovenaan ligt, met een figuur of een testwaarde. Integreer bovenste min onderste, per stuk.
\end{werkwijze}

\nujij{Bereken de oppervlakte ingesloten door $y = x^2$ en $y = 2x$.}{Snijpunten 0 en 2. $\int_0^2 (2x - x^2)\dx = 4 - \tfrac83 = \tfrac43$.}

Snijden de grafieken ook binnen het interval, dan wisselen ze van plaats en moet je daar splitsen. Zonder figuur zie je dat makkelijk over het hoofd.

\voorbeeld{}{Bereken de oppervlakte van het gebied ingesloten door de grafieken van $f(x) = x^3$ en $g(x) = x$.\label{vb:h2-wissel}}{%
\uitleg{De drie snijpunten geven twee stukken, en in elk stuk ligt een andere grafiek bovenaan. De figuur toont dat meteen; zonder figuur test je een punt in elk stuk.}
\naast[3.4cm]{\begin{stappenvrij}{6.8cm}{3.9cm}
x^3 = x \ \Rightarrow \ x = -1,\ x = 0,\ x = 1 & Drie snijpunten, dus twee stukken. \\[1ex]
\ig_{-1}^{0} (x^3 - x)\dx = \grens{\dfrac{x^4}{4} - \dfrac{x^2}{2}}{-1}{0} = \dfrac14 & Op $[-1, 0]$ ligt $x^3$ bovenaan. \\[3ex]
\ig_{0}^{1} (x - x^3)\dx = \grens{\dfrac{x^2}{2} - \dfrac{x^4}{4}}{0}{1} = \dfrac14 & Op $[0, 1]$ ligt $x$ bovenaan. \\[3ex]
A = \dfrac14 + \dfrac14 = \dfrac12 & Eén integraal over $[-1, 1]$ gaf 0.
\end{stappenvrij}}{\begin{tikzpicture}[x=1cm,y=1cm,font=\scriptsize]
  \gebiedtussen{-1}{0}{(\x)^3}{\x}
  \gebiedtussen{0}{1}{\x}{(\x)^3}
  \rooster{-1}{1}{1}{-1}{1}{1}
  \kromme{-1.1}{1.1}{(\x)^3}
  \kromme{-1.1}{1.1}{\x}
\end{tikzpicture}}}

\begin{oefeningen}
\oef{h2-tussen}
\end{oefeningen}

%========================================================================================
\opdrachtenkop

\begin{oefeningen}
\needspace{13\baselineskip}
\oef{h2-dp}

\needspace{17\baselineskip}
\oef{h2-d1}

\oef{h2-o2}

\needspace{14\baselineskip}
\oef{h2-d2}

\needspace{10\baselineskip}
\oef{h2-a1}

\oef{h2-o5}

\oef{h2-o6}

\oef{h2-o10}

\oef{h2-o11}

\oef{h2-okies}

\needspace{13\baselineskip}
\oef{h2-d3}

\oef{h2-o13}

\oef{h2-o14}

\needspace{19\baselineskip}
\oef{h2-a2}

\oef{h2-o17}

\oef{h2-o18}

\oef{h2-o19}

\oef{h2-o20}

\oef{h2-o21}

\oef{h2-o23}

\oef{h2-o24}

\oef{h2-o25}

\oef{h2-o26}

\oef{h2-o27}

\oef{h2-o28}

\oef{h2-o29}
\end{oefeningen}

%----------------------------------------------------------------------------------------
% KOPPEL: verbanden met andere leerstof
%----------------------------------------------------------------------------------------
\needspace{20\baselineskip}
\opdrachtsectie{Koppel}

\lkr{Matrices zijn gezien in het 5de jaar. Rijen en Statistiek staan volgens de README van het pakket in het 6de jaar, ná Analyse 3: kijk na of ze al aan bod kwamen. De Lorenzcurve en de Gini-coëfficiënt staan in het leerplan als context bij LPD 31.}

\begin{oefeningen}
\oef{h2-koppel-1}

\oef{h2-koppel-2}

\oef{h2-koppel-3}

\oef{h2-koppel-4}
\end{oefeningen}

%----------------------------------------------------------------------------------------
\opdrachtsectie{Test jezelf}

\begin{oefeningen}
\oef{h2-t1}

\oef{h2-t2}

\oef{h2-t5}

\oef{h2-t6}

\oef{h2-t7}

\oef{h2-t8}
\end{oefeningen}

\herhaalschema{%
  \fv{oef:h2-t1}{vb:h2-afgeleide}{oef:h2-afgeleide}%
  \fv{oef:h2-t2}{vb:h2-lin}{oef:h2-lin}%
  \fv{oef:h2-t5}{vb:h2-asopp}{oef:h2-asopp}%
  \fv{oef:h2-t6}{vb:h2-tussen}{oef:h2-tussen}%
  \fzv{oef:h2-t7}{oef:h2-o10}%
  \fzv{oef:h2-t8}{oef:h2-o19}}

\ifzelfstudie
\section*{Maak je eigen samenvatting}
Sluit dit boek en schrijf in je schrift: de twee delen van de hoofdstelling, met het bewijs in je eigen woorden, de tabel met basisprimitieven, wanneer je de stelling niet mag gebruiken, en hoe je een oppervlakte berekent met de as en tussen twee grafieken. Vergelijk daarna met de kaders.
\fi
